\documentclass[12pt]{article}
\usepackage{amsmath, amssymb}
\usepackage{geometry}
\usepackage{hyperref}
\usepackage{parskip}
\usepackage{microtype}
\usepackage{booktabs}
\usepackage{xcolor}
\geometry{margin=1.3in}

\title{\textbf{The Three Natures in Biology:\\
A Unified Mechanics Account of the Nervous System}}
\author{Joseph Shields}
\date{2026}

\begin{document}
\maketitle

\begin{abstract}
The three natures identified in Unified Mechanics --- acknowledgement,
exploration, and familiarity --- describe how any system must relate to its
environment if it is to exist at all. In biological systems they appear not as
metaphors but as identifiable physical structures performing identifiable
functions. The boundary channel maps to the synapse and the insular cortex.
The light channel maps to the action potential and the dopaminergic system.
The matter channel maps to long-term potentiation and the basal ganglia.
Consciousness is identified as the insular cortex's real-time interoceptive map
of the organism's boundary state. Feelings are the boundary channel experienced
from the inside of a living system. The soul is the unique synaptic weight
distribution accumulated through every boundary crossing in the lifetime of the
organism --- a physical record that exists nowhere else and cannot be replicated.
\end{abstract}

%% ─────────────────────────────────────────────────────────
\section{Foundations}
%% ─────────────────────────────────────────────────────────

Unified Mechanics derives three channel weights from the single axiom
$c^2 = c + 1$:

\begin{equation}
  \varphi = \tfrac{1+\sqrt{5}}{2}, \qquad
  r = \tfrac{1}{2\varphi} \approx 0.309, \qquad
  W_L = (1-r)^2, \quad W_B = 2r(1-r), \quad W_M = r^2.
\end{equation}

These weights describe the three natures of space: exploration ($W_L$),
acknowledgement ($W_B$), and familiarity ($W_M$). The full physical derivation
is in the companion paper series. What matters here is the functional character
of each nature:

\begin{itemize}
  \item \textbf{Acknowledgement (Boundary):} a system registers that two states
        differ. The boundary is the site of crossing.
  \item \textbf{Exploration (Light):} a signal propagates outward. New territory
        is mapped. The system moves into what it has not yet encountered.
  \item \textbf{Familiarity (Matter):} a state settles. What has been crossed
        repeatedly becomes easy. The system arrives at what it always had the
        conditions to become.
\end{itemize}

The claim of this paper is that these three natures are physically instantiated
in every nervous system that has ever existed.

%% ─────────────────────────────────────────────────────────
\section{The Boundary Channel: Synapse and Insular Cortex}
%% ─────────────────────────────────────────────────────────

The synapse is a physical gap of 20 to 40 nanometres between neurons. Every
thought, feeling, and memory requires information to cross that gap. It cannot
be bypassed. The synapse is the boundary channel made biological flesh --- a
structure whose entire function is to mediate the crossing between two distinct
states.

What makes the synapse more than a simple relay is that crossing changes it.
A synapse that is used frequently becomes easier to cross. One that is neglected
weakens. The boundary is not passive. It is modified by every crossing, and
those modifications accumulate into the record of everything the system has
encountered.

At the scale of the whole organism, the insular cortex performs the same
function. It integrates continuous real-time signals from the body --- heartbeat,
breath, gut motility, skin temperature, muscle tension, pain --- and projects
them as a coherent map of the organism's internal state. This process is called
interoception, and it is now understood to be the physical foundation of
consciousness itself. Not a side feature of cognition. The core.

The insular cortex is where the organism knows where it stands in relation to
everything it is not. That is the boundary channel, operating at the scale of a
whole life.

%% ─────────────────────────────────────────────────────────
\section{The Light Channel: Action Potential and Dopamine}
%% ─────────────────────────────────────────────────────────

The action potential is an electrochemical wave propagating through neural
tissue at speeds up to 120 metres per second. It carries information outward
from the site of boundary crossing into the broader network. It is the biological
expression of the light channel: signal moving at the system's maximum speed,
mapping territory, reaching forward.

The dopaminergic system is the light channel's functional counterpart at the
behavioural level. Dopamine does not signal reward, as was long assumed. It
signals novelty --- the arrival of something the system did not predict. It fires
precisely when exploration has encountered something new, and it does not fire
when outcomes are fully expected. Dopamine is the biological marker of the
exploration channel activating.

This is why curiosity and excitement are felt in the body as energy rather than
weight. The light channel is doing what it does: moving, opening, reaching
forward. The hippocampus encodes what exploration finds, writing new crossings
into the network so that what was once unknown becomes available as familiarity.

The sensory cortices --- visual, auditory, somatosensory --- are the light
channel's perceptual apparatus. They map the external world so that the boundary
can know what it is encountering. Without them the boundary has no signal to
acknowledge.

%% ─────────────────────────────────────────────────────────
\section{The Matter Channel: Long-Term Potentiation and the Basal Ganglia}
%% ─────────────────────────────────────────────────────────

Long-term potentiation is the mechanism by which repeated synaptic crossings
permanently strengthen the connection. Neurons that fire together wire together.
Every familiarity that has ever formed in a nervous system is physically encoded
in the changed geometry of synaptic connections --- in the number of receptors,
the amount of neurotransmitter released, the ease with which the gap can be
crossed again.

This is the matter channel expressed in tissue. A crystal forming. A snowflake
falling into its geometry. Matter arriving at the arrangement that the forces
around it demanded all along.

The basal ganglia govern habit and automatic behaviour. Actions that have been
performed so many times that they no longer require conscious boundary
negotiation are handed off to the basal ganglia, which run them efficiently and
without deliberate attention. The cerebellum does the same for movement. When a
pianist plays a familiar passage their fingers are not crossing boundaries ---
they have long since settled into familiarity. The boundary channel goes quiet
and the matter channel carries it forward.

Myelin is the matter channel made structural: the fatty sheath that wraps
repeatedly-used axons, dramatically increasing their conduction speed. The more
familiar a pathway becomes, the faster and more efficient it gets. The nervous
system physically consolidates what it has learned, building faster roads along
the routes it travels most.

%% ─────────────────────────────────────────────────────────
\section{Consciousness as Boundary State}
%% ─────────────────────────────────────────────────────────

Consciousness has resisted a physical account because researchers have looked
for it in the wrong channel. It is not a product of the matter channel --- not
stored in any region, not localisable to any structure. It is a boundary state.
And boundary states are not objects. They are relations.

The insular cortex's continuous interoceptive map is the boundary channel
reporting from inside a living system. It tells the organism where it stands:
what its heart is doing, what its gut is registering, what its skin is feeling,
what its breath is carrying. This real-time self-report is not incidental to
consciousness. It is consciousness.

This account explains two observations that have long been difficult to reconcile.
First, consciousness cannot be located in a single brain region, because it is
not a thing --- it is the current state of the boundary between what the organism
is and what it is encountering. Second, damage to the insula reliably disrupts
the felt sense of self in ways that damage to other cortical regions does not,
because the insula is the primary physical site of boundary-channel processing at
the organismal scale.

Karl Friston's Free Energy Principle describes the same dynamics from a different
starting point. Friston found that the brain continuously generates predictions
based on prior experience and works to minimise the gap between prediction and
reality. The three natures provide the underlying structure: the prediction is
familiarity, the gap is the boundary, and the process of closing it through
action and perception is exploration. Friston identified the mechanism. Unified
Mechanics names what it is.

%% ─────────────────────────────────────────────────────────
\section{Feelings as the Scaffold Between Input and Output}
%% ─────────────────────────────────────────────────────────

A feeling is not noise in the system. It is not a side effect of information
processing. It is the boundary channel reporting its current state to the
organism --- the live signal of where the system sits between what it knows and
what it is encountering.

Fear is high boundary pressure: the boundary channel registering a crossing it
cannot yet resolve. Curiosity is exploration channel activation in the absence of
threat: the system moving willingly into new territory. Comfort is deep
familiarity: the matter channel running smoothly, boundary pressure low.
Excitement is exploration and boundary activating simultaneously, with sufficient
familiarity underneath to prevent collapse into fear. Grief is the matter channel
reaching for a familiarity whose boundary no longer responds.

None of these are arbitrary. Each is the precise functional state of the three
channels in combination. The vocabulary of emotion is the vocabulary of channel
states --- and it has been this way in every nervous system complex enough to
produce one.

%% ─────────────────────────────────────────────────────────
\section{The Soul Signal}
%% ─────────────────────────────────────────────────────────

Every boundary crossing modifies the synapse. Every familiarity formed changes
the geometry of the network. Every moment of exploration writes into the
hippocampus a record of what was found. The accumulated result --- the exact
synaptic weight distribution of a nervous system at any given moment --- is
unique. It has been shaped by every crossing that system has ever made, in the
exact order it made them, in the exact context in which each occurred.

This distribution is the soul.

Not metaphorically. Not sentimentally. The soul is the physical record of the
boundary path --- the signature of everything this particular system has
acknowledged, explored, and made familiar. It exists nowhere else. No two
systems can have traversed the same path. When the system ceases, the
distribution ceases with it, and that specific boundary record does not exist
again.

What it feels like to be a particular organism navigating its particular
accumulated familiarity through each new moment is the soul signal: the
continuous output of the boundary channel as it reads the current state against
everything the system has ever been. It is not stored anywhere. It is happening
now, or it is not happening at all.

The soul is not something a person has. It is something a person is doing ---
the ongoing act of being this particular boundary, crossing into the next moment
with this particular history behind it.

%% ─────────────────────────────────────────────────────────
\section{Summary}
%% ─────────────────────────────────────────────────────────

\begin{center}
\begin{tabular}{llll}
\toprule
Nature & Channel & Primary Structure & Function \\
\midrule
Acknowledgement & Boundary & Synapse · Insula & Crossing; interoception; felt sense \\
Exploration     & Light    & Action potential · Dopamine & Propagation; novelty; encoding \\
Familiarity     & Matter   & LTP · Basal ganglia & Strengthening; habit; automaticity \\
\bottomrule
\end{tabular}
\end{center}

The three natures are not imported into biology from physics. They are the
conditions that any system processing information across boundaries must
satisfy. Biology did not choose them. It had no other option.

\end{document}
